\chapter{1984年全国卷（理）}

\section{单选题}

\begin{problem}

数集 \(X = \{(2n + 1)\pi\mid n\in\Z\}\) 与数集 \(Y = \{(4k \pm 1)\pi\mid k\in\Z\}\) 之间的关系是（\quad）

\choices
    {\(X \subseteq Y\)}
    {\(X \supseteq Y\)}
    {\(X = Y\)}
    {\(X \neq Y\)}
\end{problem}

\begin{answer}
C
\end{answer}

\begin{solution}
任取 \(n\in\Z\). 若 \(n=2m\)，则 \(2n+1=4m+1\)；若 \(n=2m+1\)，则 \(2n+1=4(m+1)-1\)，所以 \(X\subseteq Y\). 反过来，\(4k+1=2(2k)+1\)，且 \(4k-1=2(2k-1)+1\)，所以 \(Y\subseteq X\). 因此 \(X=Y\)，故选 C.
\end{solution}

\begin{problem}

如果圆 \(x^{2} + y^{2} + Gx + Ey + F = 0\) 与 \(x\) 轴相切于原点，那么（\quad）

\choices
    {\(F = 0\)，\(G \neq 0\)，\(E \neq 0\)}
    {\(E = 0\)，\(F = 0\)，\(G \neq 0\)}
    {\(G = 0\)，\(F = 0\)，\(E \neq 0\)}
    {\(G = 0\)，\(E = 0\)，\(F \neq 0\)}
\end{problem}

\begin{answer}
C
\end{answer}

\begin{solution}
圆与原点相切，原点必在圆上，代入方程得 \(F=0\). 此时方程化为 \(x^{2}+y^{2}+Gx+Ey=0\)，其圆心为 \(\left(-\frac{G}{2},-\frac{E}{2}\right)\). 圆在原点处与 \(x\) 轴相切，所以圆心与原点的连线垂直于 \(x\) 轴，圆心的横坐标必须为零，从而 \(G=0\). 又因为这是一个圆，半径不能为零，所以 \(E\neq0\). 因此正确选项为 C.
\end{solution}

\begin{problem}

如果 \(n\) 是正整数，那么 \(\frac{1}{8}[1 - (-1)^n](n^2 - 1)\) 的值（\quad）

\choices
    {一定是零}
    {一定是偶数}
    {是整数但不一定是偶数}
    {不一定是整数}
\end{problem}

\begin{answer}
B
\end{answer}

\begin{solution}
当 \(n\) 为偶数时，\(1-(-1)^n=0\)，原式等于零. 当 \(n\) 为奇数时，\(1-(-1)^n=2\)，且奇数的平方模 \(8\) 余 \(1\)，所以 \(n^2-1\) 是 \(8\) 的倍数，原式等于 \(\frac{n^2-1}{4}\). 进一步，\(\frac{n^2-1}{4}=\frac{n-1}{2}\cdot\frac{n+1}{2}\)，其中两个因数是相邻整数，必有一个是偶数，因此原式为偶数. 综上，原式一定是偶数，故选 B.
\end{solution}

\begin{problem}

\(\arccos (-x)\) 大于 \(\arccos x\) 的充分条件是（\quad）

\choices
    {\(x \in (0,1]\)}
    {\(x \in (-1,0)\)}
    {\(x \in [0,1]\)}
    {\(x \in \left[0, \frac{\pi}{2}\right]\)}
\end{problem}

\begin{answer}
A
\end{answer}

\begin{solution}
反余弦函数在 \([-1,1]\) 上单调递减. 若 \(x\in(0,1]\)，则 \(-x<x\)，从而 \(\arccos(-x)>\arccos x\)，所以该条件是充分的. 事实上，在反余弦有意义的范围内，\(\arccos(-x)>\arccos x\) 等价于 \(-x<x\)，即 \(x>0\). 因此正确选项为 A.
\end{solution}

\begin{problem}

如果 \(\theta\) 是第二象限角，且满足 \(\cos \frac{\theta}{2} - \sin \frac{\theta}{2} = \sqrt{1 - \sin \theta}\)，那么 \(\frac{\theta}{2}\)（\quad）

\choices
    {是第一象限角}
    {是第三象限角}
    {可能是第一象限角，也可能是第三象限角}
    {是第二象限角}
\end{problem}

\begin{answer}
B
\end{answer}

\begin{solution}
令 \(a=\frac{\theta}{2}\). 由二倍角公式，\(1-\sin\theta=1-2\sin a\cos a=(\cos a-\sin a)^2\)，所以题设等式化为 \(\cos a-\sin a=\lvert\cos a-\sin a\rvert\)，从而 \(\cos a-\sin a\geqslant0\). 由于 \(\theta\) 是第二象限角，\(a\) 的终边可能在第一象限或第三象限. 若 \(a\) 在第一象限，则 \(a\in\left(\frac{\pi}{4},\frac{\pi}{2}\right)\)，于是 \(\cos a-\sin a<0\)，不符合条件；若 \(a\) 在第三象限，则 \(a\in\left(\frac{5\pi}{4},\frac{3\pi}{2}\right)\)，于是 \(\cos a-\sin a>0\)，符合条件. 因此 \(\frac{\theta}{2}\) 是第三象限角，故选 B.
\end{solution}

\section{解答题}

\begin{problem}

已知圆柱的侧面展开图是边长为 \(2\) 与 \(4\) 的矩形，求圆柱的体积.
\end{problem}

\begin{solution}
设圆柱的底面半径为 \(r\)，高为 \(h\). 侧面展开图的两边分别是圆柱的底面周长和高，因此只有两种对应关系.

当 \(2\pi r=2\)，\(h=4\) 时，\(r=\frac{1}{\pi}\)，所以 \(V=\pi r^2h=\frac{4}{\pi}\). 当 \(2\pi r=4\)，\(h=2\) 时，\(r=\frac{2}{\pi}\)，所以 \(V=\pi r^2h=\frac{8}{\pi}\). 故圆柱的体积为 \(\frac{4}{\pi}\) 或 \(\frac{8}{\pi}\).
\end{solution}

\begin{problem}

函数 \(\log_{0.5}(x^2 + 4x + 4)\) 在什么区间上是增函数？
\end{problem}

\begin{solution}
函数的定义域由 \((x+2)^2>0\) 给出，即 \(x\neq-2\). 令 \(u=(x+2)^2\)，则当 \(x<-2\) 时，\(u\) 随 \(x\) 增大而减小；当 \(x>-2\) 时，\(u\) 随 \(x\) 增大而增大. 由于底数 \(0.5\in(0,1)\)，函数 \(\log_{0.5}u\) 是减函数，所以原函数在 \((-\infty,-2)\) 上递增，在 \((-2,+\infty)\) 上递减. 因此增函数区间为 \((-\infty,-2)\).
\end{solution}

\begin{problem}

求方程 \((\sin x + \cos x)^{2} = \frac{1}{2} \) 的解集.
\end{problem}

\begin{solution}
由 \((\sin x+\cos x)^2=\sin^2x+\cos^2x+2\sin x\cos x=1+\sin2x\)，原方程等价于 \(\sin2x=-\frac12\). 因此
\[
2x=\frac{7\pi}{6}+2n\pi\quad\text{或}\quad 2x=\frac{11\pi}{6}+2n\pi,\qquad n\in\Z.
\]
两式分别除以 \(2\)，得到 \(x=\frac{7\pi}{12}+n\pi\) 或 \(x=\frac{11\pi}{12}+n\pi\). 由于 \(\frac{11\pi}{12}+n\pi=-\frac{\pi}{12}+(n+1)\pi\)，故解集为
\[
\left\{x\mid x=\frac{7\pi}{12}+n\pi,\ n\in\Z\right\}\cup\left\{x\mid x=-\frac{\pi}{12}+n\pi,\ n\in\Z\right\}.
\]
\end{solution}

\begin{problem}

求 \(\left(|x| + \frac{1}{|x|} - 2\right)^3\) 的展开式中的常数项.
\end{problem}

\begin{solution}
令 \(t=|x|\)，则 \(t>0\)，原式为 \((t+t^{-1}-2)^3\). 利用立方公式，
\[
(t+t^{-1}-2)^3=(t+t^{-1})^3-6(t+t^{-1})^2+12(t+t^{-1})-8.
\]
其中 \((t+t^{-1})^3=t^3+3t+3t^{-1}+t^{-3}\) 不含常数项，\((t+t^{-1})^2=t^2+2+t^{-2}\) 的常数项为 \(2\)，其余的 \(12(t+t^{-1})\) 不含常数项. 因此原式的常数项为 \(-6\times2-8=-20\).
\end{solution}

\begin{problem}

求 \(\displaystyle\lim_{n\to \infty}\frac{1 - 2^n}{3^n + 1}\) 的值.
\end{problem}

\begin{solution}
分子分母同除以 \(3^n\)，得
\[
\frac{1-2^n}{3^n+1}=\frac{3^{-n}-\left(\frac{2}{3}\right)^n}{1+3^{-n}}.
\]
当 \(n\to\infty\) 时，\(3^{-n}\to0\)，\(\left(\frac23\right)^n\to0\)，所以极限为 \(0\).
\end{solution}

\begin{problem}

要排一张有 \(6\) 个歌唱节目和 \(4\) 个舞蹈节目的演出节目单，任何两个舞蹈节目不得相邻，问有多少种不同的排法.
\end{problem}

\begin{solution}
先排 \(6\) 个不同的歌唱节目，有 \(6!\) 种排法. 歌唱节目排好后形成 \(7\) 个空位：两个端点空位以及相邻歌唱节目之间的 \(5\) 个空位. 为使舞蹈节目互不相邻，至多在每个空位放置一个舞蹈节目，因此从 \(7\) 个空位中选出 \(4\) 个，有 \(\mathrm C_7^4\) 种选法. 再排列 \(4\) 个不同的舞蹈节目，有 \(4!\) 种排法. 所以总排法数为 \(\mathrm C_7^4\cdot4!\cdot6!=\mathrm A_7^4\cdot6!\).
\end{solution}

\begin{problem}

设 \(H(x) = \begin{cases}
0, & x \leqslant 0\\
1, & x > 0
\end{cases}\)，画出函数 \(y = H(x - 1)\) 的图象.
\end{problem}

\begin{solution}
令 \(u=x-1\). 当 \(x\leqslant1\) 时，\(u\leqslant0\)，所以 \(H(x-1)=0\)；当 \(x>1\) 时，\(u>0\)，所以 \(H(x-1)=1\). 因此
\[
y=H(x-1)=\begin{cases}
0,&x\leqslant1\\
1,&x>1
\end{cases}
\]
图象由直线 \(y=0\) 上满足 \(x\leqslant1\) 的部分和直线 \(y=1\) 上满足 \(x>1\) 的部分组成，点 \((1,0)\) 取到，点 \((1,1)\) 不取到.
\end{solution}

\begin{problem}

画出极坐标方程 \((\rho - 2)\left(\theta - \frac{\pi}{4}\right) = 0\)（\(\rho > 0\)）的曲线.
\end{problem}

\begin{solution}
因为两个因式的乘积为零，所以必须有 \(\rho=2\) 或 \(\theta=\frac{\pi}{4}\). 当 \(\rho=2\) 时，极角任意，得到以极点为圆心、半径为 \(2\) 的圆；当 \(\theta=\frac{\pi}{4}\) 且 \(\rho>0\) 时，得到沿与极轴正向夹角为 \(\frac{\pi}{4}\) 的方向且不含极点的半直线. 因此所求曲线是这个圆与该半直线的并集，二者交于极坐标点 \(\left(2,\frac{\pi}{4}\right)\).
\end{solution}

\begin{problem}

已知三个平面两两相交，有三条交线. 求证：这三条交线交于一点或互相平行.
\end{problem}

\begin{solution}
设三个平面为 \(\alpha\)，\(\beta\)，\(\gamma\)，并记 \(l_1=\alpha\cap\beta\)，\(l_2=\beta\cap\gamma\)，\(l_3=\gamma\cap\alpha\). 直线 \(l_1\) 与 \(l_2\) 都在平面 \(\beta\) 内，因此它们或者相交，或者平行.

\begin{enumerate}
\item 若 \(l_1\) 与 \(l_2\) 相交于点 \(P\)，则 \(P\in\alpha\) 且 \(P\in\gamma\)，所以 \(P\in\alpha\cap\gamma=l_3\). 因此三条交线交于一点 \(P\).
\item 若 \(l_1\parallel l_2\)，则 \(l_1\) 与 \(l_3\) 不可能相交. 否则交点 \(P\) 同时属于 \(\alpha\)，\(\beta\) 和 \(\gamma\)，从而 \(P\in l_2\)，这与 \(l_1\parallel l_2\) 矛盾. 因此，\(l_1\) 与 \(l_3\) 是平面 \(\alpha\) 内的两条平行直线，即 \(l_1\parallel l_3\). 同理，在平面 \(\gamma\) 内有 \(l_2\parallel l_3\)，所以三条交线互相平行.
\end{enumerate}
\end{solution}

\begin{problem}

设 \(c\)，\(d\)，\(x\) 为实数，\(c \neq 0\)，\(x\) 为未知数. 讨论方程 \(\log_{\left(cx + \frac{d}{x}\right)} x = -1\) 在什么情况下有解，有解时求出它的解.
\end{problem}

\begin{solution}
由对数的定义，原方程有意义必须满足 \(x>0\)，\(cx+\frac{d}{x}>0\)，且 \(cx+\frac{d}{x}\neq1\). 由 \(\log_{\left(cx+\frac{d}{x}\right)}x=-1\)，得
\[
x=\left(cx+\frac{d}{x}\right)^{-1},
\]
所以
\[
cx^2+d=1.
\]
由于 \(c\neq0\)，唯一可能的正数解为 \(x=\sqrt{\frac{1-d}{c}}\). 要使这个数为正数，必须有 \(\frac{1-d}{c}>0\)，即分两种情况：\(c>0\) 且 \(d<1\)，或 \(c<0\) 且 \(d>1\).

还需排除 \(x=1\)，因为对数的底数不能等于 \(1\). 由 \(x^2=\frac{1-d}{c}\) 可知 \(x=1\) 等价于 \(c=1-d\). 因此，在 \(c>0\)，\(d<1\)，\(c\neq1-d\)，或 \(c<0\)，\(d>1\)，\(c\neq1-d\) 时，令 \(x=\sqrt{\frac{1-d}{c}}\)，便有 \(cx+\frac{d}{x}=\frac{1}{x}>0\)，且底数不等于 \(1\)，从而确实满足原方程. 故有解时的解为 \(x=\sqrt{\frac{1-d}{c}}\).
\end{solution}

\begin{problem}

设 \(p \neq 0\)，实系数一元二次方程 \(z^2 - 2pz + q = 0\) 有两个虚数根 \(z_1\)，\(z_2\). 再设 \(z_1\)，\(z_2\) 在复平面内的对应点是 \(Z_1\)，\(Z_2\).

求以 \(Z_1\)，\(Z_2\) 为焦点且经过原点的椭圆的长轴的长.
\end{problem}

\begin{solution}
因为方程的系数为实数，且两个根为非实数，所以两根互为共轭复数. 设 \(z_1=p+\i r\)，\(z_2=p-\i r\)，其中 \(r\neq0\). 由根与系数的关系，
\[
q=z_1z_2=(p+\i r)(p-\i r)=p^2+r^2>0.
\]
原点在椭圆上，而焦点对应的复数为 \(z_1,z_2\)，所以原点到两个焦点的距离分别为 \(\lvert z_1\rvert=\sqrt{q}\) 和 \(\lvert z_2\rvert=\sqrt{q}\). 椭圆上任一点到两焦点距离之和等于长轴长，因此所求长轴的长为 \(\sqrt q+\sqrt q=2\sqrt q\).
\end{solution}

\begin{problem}

求经过定点 \(M\left(1,2\right)\)，以 \(y\) 轴为准线，离心率为 \(\frac{1}{2}\) 的椭圆的左顶点的轨迹方程.
\end{problem}

\begin{solution}
因为定点 \(M(1,2)\) 的横坐标为正，而椭圆的离心率小于 \(1\)，若 \(y\) 轴是右准线，则整个椭圆都在 \(y\) 轴左侧，不可能经过 \(M\). 所以 \(y\) 轴是左准线，椭圆的长轴平行于 \(x\) 轴.

设左顶点为 \(A(x,y)\)，椭圆的半长轴为 \(a_0\)，焦半距为 \(c_0\). 由离心率为 \(\frac12\)，得 \(c_0=\frac{a_0}{2}\). 左准线的方程为 \(X=h-\frac{a_0}{e}=0\)，其中 \((h,y)\) 是椭圆中心，所以 \(h=2a_0\). 左顶点的横坐标为 \(x=h-a_0=a_0\)，左焦点 \(F\) 的坐标为 \(\left(h-c_0,y\right)=\left(\frac{3x}{2},y\right)\).

由椭圆的定义，点 \(M\) 到焦点的距离等于离心率乘以点 \(M\) 到准线的距离. 点 \(M(1,2)\) 到 \(y\) 轴的距离为 \(1\)，所以 \(MF=\frac12\). 由两点间距离公式，
\[
\left(\frac{3x}{2}-1\right)^2+(y-2)^2=\frac14.
\]
整理得
\[
9\left(x-\frac23\right)^2+4(y-2)^2=1.
\]
该椭圆的横坐标范围为 \(\left[\frac13,1\right]\)，符合左顶点位于准线右侧的条件，所以所求轨迹方程为 \(9\left(x-\frac23\right)^2+4(y-2)^2=1\).
\end{solution}

\begin{problem}

在 \(\triangle ABC\) 中，\(\angle A\)，\(\angle B\)，\(\angle C\) 所对的边分别为 \(a\)，\(b\)，\(c\)，且 \(c = 10\)，\(\frac{\cos A}{\cos B} = \frac{b}{a} = \frac{4}{3}\)，\(P\) 为 \(\triangle ABC\) 的内切圆上的动点. 求点 \(P\) 到顶点 \(A\)，\(B\)，\(C\) 的距离的平方和的最大值与最小值.
\end{problem}

\begin{solution}
由正弦定理，\(\frac{b}{a}=\frac{\sin B}{\sin A}\). 结合题设 \(\frac{\cos A}{\cos B}=\frac{\sin B}{\sin A}\)，得 \(\sin A\cos A=\sin B\cos B\)，即 \(\sin2A=\sin2B\). 由于 \(a\neq b\)，所以 \(A\neq B\)，只能有 \(A+B=\frac\pi2\)，从而 \(C=\frac\pi2\). 因此 \(c\) 是斜边，且 \(a:b=3:4\). 由勾股定理，\(a=6\)，\(b=8\)，\(c=10\).

三角形面积为 \(\frac12ab=24\)，半周长为 \(s=\frac{6+8+10}{2}=12\)，所以内切圆半径为 \(r=\frac{24}{12}=2\). 建立直角坐标系，取 \(C=(0,0)\)，\(A=(8,0)\)，\(B=(0,6)\). 则内心为 \((2,2)\)，内切圆方程为 \((x-2)^2+(y-2)^2=4\). 设动点 \(P=(x,y)\)，则
\[
S=PA^2+PB^2+PC^2=(x-8)^2+y^2+x^2+(y-6)^2+x^2+y^2=3x^2+3y^2-16x-12y+100.
\]
由内切圆方程得 \(x^2+y^2=4x+4y-4\)，代入上式可得 \(S=88-4x\). 又因为圆心横坐标为 \(2\)，半径为 \(2\)，所以 \(0\leqslant x\leqslant4\). 于是 \(S\) 的最大值为 \(88\)，最小值为 \(72\)，且分别在 \(x=0\) 和 \(x=4\) 时取得.
\end{solution}

\begin{problem}

设 \(a > 2\)，给定数列 \(\{x_n\}\)，其中 \(x_1 = a\)，\(x_{n+1} = \frac{x_n^2}{2(x_n - 1)}\)（\(n = 1\)，\(2\)，\(\cdots\)）. 求证：

\begin{enumerate}
\item \(x_n > 2\)，且 \(\frac{x_{n + 1}}{x_n} < 1\)（\(n = 1\)，\(2\)，\(\cdots\)）；
\item 如果 \(a \leqslant 3\)，那么 \(x_n \leqslant 2 + \frac{1}{2^{n - 1}}\)（\(n = 1\)，\(2\)，\(\cdots\)）；
\item 如果 \(a > 3\)，那么当 \(n \geqslant \frac{\lg \frac{a}{3}}{\lg \frac{4}{3}}\) 时，必有 \(x_{n + 1} < 3\).
\end{enumerate}
\end{problem}

\begin{solution}
由递推式，
\[
x_{n+1}-2=\frac{x_n^2}{2(x_n-1)}-2=\frac{(x_n-2)^2}{2(x_n-1)}.
\]

\begin{enumerate}
\item 当 \(n=1\) 时，\(x_1=a>2\). 若 \(x_n>2\)，则上式右端为正，且分母不为零，从而 \(x_{n+1}>2\). 由数学归纳法，\(x_n>2\) 对所有正整数 \(n\) 成立. 又因为
\[
\frac{x_{n+1}}{x_n}=\frac{x_n}{2(x_n-1)}<1
\]
等价于 \(x_n>2\)，所以 \(\frac{x_{n+1}}{x_n}<1\) 对所有正整数 \(n\) 成立.
\item 由第（1）问知 \(x_n>2\)，因此 \(0<\frac{x_n-2}{x_n-1}<1\)，从而
\[
x_{n+1}-2=\frac{x_n-2}{2}\cdot\frac{x_n-2}{x_n-1}<\frac{x_n-2}{2}.
\]
当 \(a\leqslant3\) 时，\(x_1-2=a-2\leqslant1\). 反复使用上式，得到
\[
x_n-2\leqslant\frac{a-2}{2^{n-1}}\leqslant\frac{1}{2^{n-1}},
\]
即 \(x_n\leqslant2+\frac{1}{2^{n-1}}\).
\item 对任意满足 \(x_k>3\) 的正整数 \(k\)，由第（1）问中的比值公式有
\[
\frac{x_{k+1}}{x_k}=\frac{x_k}{2(x_k-1)}<\frac34.
\]
记 \(L=\frac{\lg\frac{a}{3}}{\lg\frac43}\). 任取正整数 \(n\geqslant L\). 若 \(x_n\leqslant3\)，则由 \(\frac{x_{n+1}}{x_n}<1\) 可得 \(x_{n+1}<x_n\leqslant3\). 若 \(x_n>3\)，则由第（1）问知 \(x_1>x_2>\cdots>x_n>3\)，于是
\[
x_{n+1}=a\prod_{k=1}^{n}\frac{x_{k+1}}{x_k}<a\left(\frac34\right)^n\leqslant a\left(\frac34\right)^L=3.
\]
最后一个等式是因为 \(\left(\frac43\right)^L=\frac{a}{3}\). 两种情形均得到 \(x_{n+1}<3\)，命题得证.
\end{enumerate}
\end{solution}

\section{附加题}

\begin{problem}

如图，已知圆心为 \(O\)，半径为 \(1\) 的圆与直线 \(l\) 相切于点 \(A\)，一动点 \(P\) 自切点 \(A\) 沿直线 \(l\) 向右移动时，取 \(\widehat{AC}\) 的长为 \(\frac{2}{3} AP\)，直线 \(PC\) 与直线 \(AO\) 交于点 \(M\). 又知当 \(AP = \frac{3\pi}{4}\) 时，点 \(P\) 的速度为 \(v\)，求这时点 \(M\) 的速度.

\bitmapfigure[float=inline,max width=0.42\linewidth]{img/1984/national_science/q20_fig1.png}
\FloatBarrier
\end{problem}

\begin{solution}
设 \(\theta=\angle AOC\)（用弧度表示），并设 \(x=AP\)，\(y=AM\). 过 \(C\) 作 \(CD\perp AO\)，则 \(DC=\sin\theta\)，\(AD=1-\cos\theta\).
由于圆的半径为 \(1\)，弧长 \(\widehat{AC}=\theta\). 由题设 \(\widehat{AC}=\frac{2}{3}AP=\frac{2x}{3}\)，故 \(\theta=\frac{2x}{3}\).
直角三角形 \(\triangle APM\) 与 \(\triangle DCM\) 相似，因此
\[
\frac{AM}{AP}=\frac{DM}{DC}.
\]
代入 \(AM=y\)，\(AP=x\)，\(DM=y-(1-\cos\theta)\)，以及 \(DC=\sin\theta\)，得到
\[
\frac{y}{x}=\frac{y-(1-\cos\theta)}{\sin\theta}.
\]
解得
\[
y=\frac{x(1-\cos\theta)}{x-\sin\theta}.
\]
对 \(x\) 求导，
\[
\frac{dy}{dx}=\frac{\left(1-\cos\theta+\frac{2x}{3}\sin\theta\right)(x-\sin\theta)-x(1-\cos\theta)\left(1-\frac{2}{3}\cos\theta\right)}{(x-\sin\theta)^2}.
\]
当 \(x=\frac{3\pi}{4}\) 时，\(\theta=\frac\pi2\)，所以 \(\sin\theta=1\)，\(\cos\theta=0\)，从而
\[
\left.\frac{dy}{dx}\right|_{x=\frac{3\pi}{4}}=\frac{\left(1+\frac\pi2\right)\left(\frac{3\pi}{4}-1\right)-\frac{3\pi}{4}}{\left(\frac{3\pi}{4}-1\right)^2}=\frac{2(3\pi^2-4\pi-8)}{(3\pi-4)^2}.
\]
由 \(x=AP\) 且点 \(P\) 的速度为 \(v\)，有 \(\frac{dx}{dt}=v\)，所以
\[
\frac{dy}{dt}=\frac{dy}{dx}\frac{dx}{dt}=\frac{2(3\pi^2-4\pi-8)}{(3\pi-4)^2}v.
\]
因此这时点 \(M\) 的速度为 \(\frac{2(3\pi^2-4\pi-8)}{(3\pi-4)^2}v\).
\end{solution}
